\chapter{Simulation Configuration Files}
\label{app:config}

\section{Complete \texttt{omnetpp.ini}}
\label{app:omnetpp_ini}

The complete \texttt{omnetpp.ini} configuration file used in all simulations is listed
below. Comments explain the purpose of each parameter group.

\begin{lstlisting}[caption={Complete \texttt{omnetpp.ini} simulation configuration.},
                   label={lst:omnetpp_ini}]
[General]
network = inet.examples.thesis_wmn.MeshNetwork
sim-time-limit = 50s
record-eventlog = false

# 1. DISABLE STATIC ROUTES (force AODV/DSDV to handle routing)
*.configurator.addStaticRoutes = false

# 2. NODE COUNT SWEEP
*.numNodes = ${N=10,20,30,40,50}

# 3. MOBILITY: Stationary (MassMobility with speed=0)
# Using MassMobility instead of StationaryMobility to avoid
# a macOS Qtenv crash in OMNeT++ 6.3 on Apple Silicon.
**.mobility.bounds = "1000m 1000m 0m"
**.mobility.typename        = "MassMobility"
**.mobility.initialX        = uniform(0m, 1000m)
**.mobility.initialY        = uniform(0m, 1000m)
**.mobility.initialZ        = 0m
**.mobility.speed           = 0mps
**.mobility.changeInterval  = 1s
**.mobility.angleDelta      = 0deg
**.mobility.initFromDisplayString = false

# Show radio range circles in the Qtenv GUI
**.wlan[*].radio.displayCommunicationRange = true

# 4. IEEE 802.11g RADIO MEDIUM
**.radioMedium.typename = "Ieee80211ScalarRadioMedium"
**.radioMedium.backgroundNoise.power = -90dBm
**.wlan[*].typename       = "Ieee80211Interface"
**.wlan[*].radio.typename = "Ieee80211ScalarRadio"
**.wlan[*].mac.typename   = "Ieee80211Mac"
**.wlan[*].mac.opMode     = "g(mixed)"
**.wlan[*].mac.bitrate    = 54Mbps
**.wlan[*].radio.transmitter.communicationRange = 100m

# 5. ENERGY MODEL
**.energyStorage.typename                           = "IdealEpEnergyStorage"
**.energyStorage.initialCapacity                    = 100J
**.wlan[*].radio.energyConsumer.typename            = "StateBasedEpEnergyConsumer"
**.wlan[*].radio.energyConsumer.transmitterPowerConsumption = 1.65W
**.wlan[*].radio.energyConsumer.receiverPowerConsumption    = 1.4W
**.wlan[*].radio.energyConsumer.sleepPowerConsumption       = 0.1W

# 6. TRAFFIC: Many-to-one UDP (all nodes -> host[0])
**.host[*].numApps = 1

# All nodes send to the sink
**.host[*].app[0].typename       = "UdpBasicApp"
**.host[*].app[0].destAddresses  = "host[0]"
**.host[*].app[0].destPort       = 5000
**.host[*].app[0].messageLength  = 1000B
**.host[*].app[0].sendInterval   = uniform(0.5s, 1s)
**.host[*].app[0].startTime      = uniform(1s, 5s)
**.host[*].app[0].stopTime       = 49s

# Sink receives on port 5000
**.host[0].app[0].typename   = "UdpSink"
**.host[0].app[0].localPort  = 5000

# Enable end-to-end delay recording at the sink
**.host*.app[0].endToEndDelay.result-recording-modes = +mean
**.host*.app[0].endToEndDelay.scalar-recording = true

# --- PROTOCOL CONFIGURATIONS ---

[Config AODV_Test]
description = "Simulate Reactive AODV (RFC 3561)"
**.routingApp.typename = "Aodv"

[Config DSDV_Test]
description = "Simulate Proactive DSDV (Perkins & Bhagwat 1994)"
**.routing.typename = "Dsdv"
\end{lstlisting}

\section{Complete \texttt{MeshNetwork.ned}}
\label{app:ned}

\begin{lstlisting}[caption={\texttt{MeshNetwork.ned} --- network definition file.},
                   label={lst:ned}]
package inet.examples.thesis_wmn;

import inet.networklayer.configurator.ipv4.Ipv4NetworkConfigurator;
import inet.node.inet.ManetRouter;
import inet.physicallayer.wireless.common.medium.RadioMedium;

network MeshNetwork
{
    parameters:
        int numNodes;
        @display("bgb=1000,1000"); // 1000m x 1000m deployment area

    submodules:
        configurator: Ipv4NetworkConfigurator {
            @display("p=100,100");
        }
        radioMedium: RadioMedium {
            @display("p=100,200");
        }
        host[numNodes]: ManetRouter {
            @display("p=500,500");
        }
}
\end{lstlisting}
